\section{Introduction}

Fuel poverty research has made household energy hardship visible, but the recent politics of domestic energy requires a further question. A household may not be classified as fuel poor at one point in time and yet remain highly exposed to the next fossil-fuel price shock; equally, a household may receive a low-carbon technology without receiving meaningful protection from volatile retail energy costs. This paper therefore reframes household energy inequality from a static problem of current fuel poverty into a dynamic problem of uneven exposure to future fossil-fuel price volatility.
% Citation TODO: fuel poverty metrics, energy vulnerability, energy price resilience.

This reframing is significant for three reasons. First, energy-price shocks are not merely temporary deviations from a normal affordability problem. They reveal which households have high required energy needs, thin residual incomes, limited ability to absorb bill increases, and little control over the building or tariff conditions that shape their exposure. Second, actual expenditure can understate hardship when households respond to high prices by underheating, self-rationing, or delaying payment. A risk framework must therefore start from required energy services and then ask how costs change under plausible price and policy scenarios. Third, domestic net zero is increasingly presented as a route to lower bills and improved resilience, but the household-protection effect is not automatic. It depends on whether technical change reduces required demand, changes the household's fuel-price exposure, and is accessible under the household's tenure, dwelling form, tariff options, and local governance constraints.
% Citation TODO: underheating/required expenditure; retrofit and net-zero policy claims; tariff and price-shock literature.

Existing studies provide important building blocks but do not fully answer this paper's question. Fuel poverty indicators identify current need and support policy targeting, while energy vulnerability research shows that hardship is produced through income, housing quality, health, tenure relations, social support, and institutional access. Recent work on energy price resilience, energy poverty premiums, underheating, and crisis-period demand reduction moves closer to household price-risk analysis. Retrofit and low-carbon technology studies also show that fabric improvements, heat pumps, PV, batteries, and flexibility can alter energy demand and bills. Yet these literatures often remain separated: fuel poverty metrics are usually current-condition indicators; retrofit studies often focus on average energy or bill effects; low-carbon adoption research examines who installs technologies more than who receives price-risk protection; and price-shock studies rarely connect future volatility to local housing governance and intervention feasibility at property or small-area scale.
% Citation TODO: fuel poverty indicators; energy vulnerability; household energy price resilience; low-carbon adoption inequality; retrofit evaluation.

The gap is especially important in the United Kingdom because decarbonisation and de-risking can diverge. Electrification can reduce direct household gas use and emissions, but it may not reduce household fuel price risk if electricity prices remain high relative to gas or continue to transmit gas-market shocks through wholesale and retail bill structures. Similarly, PV, batteries, and smart tariffs may hedge electricity-price exposure for some households while excluding others because of roof access, leasehold governance, smart-meter access, digital capability, load flexibility, or capital cost. The policy question is therefore not whether a measure is low carbon in engineering terms, but whether it reduces upper-tail household burden risk for the households most exposed to price volatility.
% Citation TODO: UK electricity-gas price ratio, gas pass-through, smart tariff equity, PV/battery access.

Westminster is an analytically revealing case rather than a merely convenient location. High average wealth can obscure fuel-price vulnerability among households facing high housing costs, inefficient or hard-to-retrofit dwellings, private rental constraints, leasehold governance, heritage restrictions, and limited access to low-carbon assets. Dense flats, social housing blocks, private rented weak-fabric dwellings, electric-heated inefficient flats, large inefficient gas-heated homes, and conservation-area properties create different routes into exposure and different routes out of it. These conditions make Westminster a strong test of whether domestic net zero can operate as household price-risk protection, and for whom.
% Citation TODO: Westminster/London housing, retrofit, PRS, heritage, fuel poverty evidence.

The paper's starting point is therefore not that domestic net zero measures are inherently protective. A dwelling can become lower carbon without becoming less exposed to household price risk if the intervention leaves required heat demand high, shifts the household onto a more expensive fuel, excludes the occupant from favourable tariffs, or depends on building-level decisions outside the household's control. Conversely, the same intervention may be protective where fabric improvement, heating-system efficiency, tariff design, finance, and governance align. The key distributional question is who receives this protection and who remains exposed.

The central research question is: to what extent, and under what housing, market, and governance conditions, do domestic net zero interventions reduce unequal household exposure to fossil-fuel price volatility in Westminster?

The paper addresses three linked questions:
\begin{enumerate}[label=\arabic*.]
    \item Which household-dwelling types face the largest required energy burden increase under fuel price shocks?
    \item How do income, housing fabric, heating system, tenure, and vulnerability jointly produce fuel price risk?
    \item How much risk reduction do retrofit, clean heat, PV, battery storage, smart tariffs, and integrated packages deliver, and for whom?
\end{enumerate}

The contribution is fourfold. Conceptually, the paper distinguishes fuel poverty as current hardship, energy burden as a cost-to-income relation, and household fuel price risk as future-facing exposure, sensitivity, and limited adaptive capacity under price-shock scenarios. Methodologically, it develops a transparent income--cost and scenario framework that links required energy demand, local residual-income distributions, fuel-price shocks, intervention feasibility, and upper-tail burden risk. Empirically, it applies this framework to Westminster's heterogeneous housing stock, where high-cost urban housing, tenure governance, flats, heritage constraints, and social housing delivery pathways are central mechanisms rather than contextual details. For policy, it evaluates domestic net zero interventions as conditional de-risking tools, not as automatic bill-reduction or emissions-reduction proxies.

The modelling strategy follows directly from this argument. The baseline layer estimates required energy burden and the probability of crossing a burden threshold, because these measures identify where current cost and income conditions make exposure plausible even without observing each occupant's income. The price-shock layer then varies fuel prices, standing charges, tariff structures, and mitigation to estimate how exposure changes under future conditions. The intervention layer finally asks whether fabric retrofit, clean heat, PV, storage, smart tariffs, or integrated packages reduce upper-tail burden risk once access constraints are included. The equations are therefore not a detached technical apparatus; they are the operational form of the paper's central claim that current hardship, future exposure, and conditional protection must be measured separately before their distribution can be compared.
